% PDF 32--58；原书第 1--27 页。
\chapter{大语言模型导论}
\label{ch:intro}
\chaptercredit{1}
\chapterabstract{本章先讨论自然语言处理的历史背景与演进。从人类语言能力的起源出发，说明人类语言如何逐步成为计算建模的对象，并强调生物演化与技术之间复杂联系的重要性。本章沿时间脉络梳理计算领域中语言模型的发展，展示早期简单模型如何演变为今天熟悉的大语言模型，讨论促成转变的关键因素，包括算法进展、算力和数据可获得性。随后界定大语言模型，帮助读者理解它们在当代人工智能范式中的意义。最后概述后续各章，使读者了解全书主题的广度与深度。}

% PDF 32--33。
\section{引言}
\label{sec:intro-opening}
2022 年 11 月，总部位于旧金山的科技公司 OpenAI 宣布公开发布 ChatGPT。这是一款基于网页、经训练能够回答用户查询的聊天机器人。随之而来的报道与传播在全球引起轰动。到 2023 年 1 月，该平台的月活跃用户已经超过 1 亿，使它成为有史以来增长最快的消费应用。人们围绕 ChatGPT 展开讨论，关注它生成的文本如何近似人类表达、它回答各学科技术问题时的深度、它在 GRE 和 LSAT 等标准化考试中的出色表现，以及它限制回答争议性问题的安全约束（safety guardrails）。随后数月，市场上出现多款竞品，包括 Microsoft 的 Bing Chat（2023 年 2 月）、Google 的 Bard（2023 年 3 月）、Anthropic 的 Claude（2023 年 3 月）和百度（Baidu）的文心一言（Ernie，2023 年 8 月）。

几乎与此同时，研究者在另一类生成模型（generative model）上取得迅速进展，即文本到图像模型（text-to-image model）。这类模型接收用户写下的描述，生成与描述相符的数字图像。自 2022 年起，当时最先进的（state of the art，SOTA）文本到图像模型就能针对仿佛无所不包的提示主题生成接近照片的图像，一些广为流传的例子甚至骗过不少人（\refcite{ch01-di-placido}{Di Placido, 2023}）。这些程序还可以为任意题材作图，模仿著名艺术家的风格以及历代艺术流派，令人印象深刻。参与者包括 OpenAI（DALL-E 2，2022 年 4 月）、Midjourney（2022 年 7 月）和 StabilityAI（Stable Diffusion，2022 年 8 月）。2022 和 2023 年引发热议的相关人工智能应用，还包括合成模仿名人说话和歌唱的声音，以及篡改实时视频、改变人物外貌或讲话内容的深度伪造（deep fakes）。

这些程序在细节上各异，却有共同的技术基础，也就是大语言模型。这项技术迅速进入公众视野，带动了对生成式人工智能公司的投资热潮。到 2024 年 2 月，OpenAI 估值飙升至 800 亿美元，超过 2021 年的五倍。NVIDIA 是人工智能技术所需图形处理器（Graphics Processing Unit，GPU）的主要制造商。2022 年 10 月之后约一年半内，其市值增至十倍以上，最终于 2024 年 6 月超过 Microsoft，至少一度成为全球市值最高的公司。全球生成式人工智能市场的估值从 2021 年的 82 亿美元跃升至 2022 年底的 290 亿美元，预计到 2030 年将超过 6\,670 亿美元（\refcite{ch01-fortune}{Fortune, 2023}；\refcite{ch01-jin}{Jin and Kruppa, 2023}；\refcite{ch01-valuates}{Valuates, 2023}）。金融市场和私人投资者预期，未来十年这一领域将大幅增长。

人工智能投资的增长也带来了公众审视。围绕基于大语言模型的人工智能，人们经常讨论社会、政治和生存风险，以及经济影响、伦理与长期后果。艺术家担心人工智能生成的艺术作品和文字侵入他们的领域。教育工作者思考，在有 ChatGPT 的世界里，大学入学申请文章是否还能被视为申请人亲笔所写。世界各地的政府都在考虑监管人工智能研究，使用从互联网上抓取的数据训练大语言模型是否合法，也在多起诉讼中接受审理。短短两年，人工智能的公共讨论已在各个方面发生巨变。熟悉近期进展的人工智能专家可能对这些平台的成就并不惊讶，但大众由此第一次领略到过去十年机器学习领域发生的变革。2022 年可说是大语言模型进入全球公众视野的起点。

这些大语言模型是什么？它们为何具备如此惊人的能力？它们能够精细吸收、处理并运用自然语言，背后有什么机制？过去几十年的机器学习研究，又如何让大语言模型得以登上世界舞台？要回答这些问题，可以先退一步，考察这些能力所依赖的基础系统，即人类语言。

% PDF 34--36。
\section{自然语言}
在人类以外的动物交流方式中，找不到与自然语言相同的形式。交流是个体之间交换思想、感受、观念或信息的过程。没有其他交流方式比人类自然语言更具表现力、更高效，也更抽象。从根本上说，语言通过以复杂方式组合的符号、手势和声音来精确编码意义。人类能准确传达有关世界的意义，这被认为是人类取得地球生态优势的最重要因素之一。有了语言，意义可以从一个个体传递给许多人，使有用的信息快速传播，惠及收到信息的人。

不难想象，自然语言交流会怎样改变早期人类群体的生活。人类起源于群居祖先，而交流对群居物种的演化成就必不可少。不妨设想，生活在一个没有语言交流的人类群体中。缺乏自然语言时，人们能以多高的效率解决纠纷、共享有利于大家的资源位置，或者告知群体威胁来自哪个方向？这些事并非做不到，地球上其他物种也使用表达能力弱于自然语言的方式交流。不过，人类适应环境的潜力将受到很大限制。无论结果更好还是更坏，我们都会成为与今天很不相同的生物。

语言拓展了我们在行为上适应环境的方式。一个物种是被限制在固定的生活史（life history）中，成为生态特化种（ecological specialist），还是拥有灵活的生活史，成为生态广适种（ecological generalist），这一差别至关重要。以鲑科（Salmonidae）为例，其中包括多种鲑鱼和鳟鱼。鲑科有独特的生活史。幼鱼通常在出生的河流和溪流中成长，达到一定年龄和体型后，迁徙极远的距离，到资源更丰富的环境摄食，直至长成成鱼。达到繁殖年龄时，成鱼返回出生水域产卵，从而开启下一代的循环。这一生活史在鲑科个体中几乎固定不变，因为对鲑鱼或鳟鱼而言，采取不同于亲代的生活方式代价可能很高。

相比之下，人类的生活史多种多样，而个体的生存能力大体不受影响。自然语言如何赋予我们这一能力？道理出奇简单。我们能够精确地传达意义，因而可以了解世界上的事，而不必等到亲历信息所描述的事件或事实。即使在时间或空间上远离这些事件，我们也能间接获知。这是演化带给人类的一项非凡能力。

\begin{figure}[!htbp]
\centering
% 重绘自原书第 4 页的图 1.1。
\begin{tikzpicture}[
  font=\small\sffamily,
  entity/.style={circle,minimum size=12mm,inner sep=0pt,fill=blue!53!cyan,text=white,align=center},
  edge/.style={-{Stealth[length=2mm]},draw=black!62,thick},
  relation/.style={font=\footnotesize\sffamily,fill=white,inner sep=1.2pt,align=center}]
\node[entity] (tiger) at (2.6,3.1) {老虎};
\node[entity] (predator) at (4.25,4.35) {捕食者};
\node[entity] (human) at (7.5,4.55) {人};
\node[entity,fill=orange!75!yellow,text=black] (self) at (7.7,1.85) {自身};
\node[entity] (bush) at (1.2,1.1) {响动的\\灌木丛};
\node[entity] (wind) at (4.25,.45) {风};
\node[entity,fill=black!13,draw=black!55,densely dotted,text=black] (america) at (10.35,5.4) {北美洲};
\draw[edge] (tiger.north east) to[bend left=25] node[relation,below right] {属于} (predator.west);
\draw[edge] (predator.east) to[bend right=10] node[relation,above] {有些会吃} (human.west);
\draw[edge] (self.north) to[bend left=13] node[relation,right] {是} (human.south);
\draw[edge] (tiger.east) to[out=0,in=180] node[relation,pos=.43,above] {构成严重威胁} (self.west);
\draw[edge] (tiger.west) to[out=180,in=90] node[relation,near start,left] {有时藏在其中\\并引起响动} (bush.north);
\draw[edge] (bush.north east) to[out=24,in=180] node[relation,pos=.43,above] {可能构成威胁} (self.west);
\draw[edge] (wind.west) to[out=170,in=-12] node[relation,pos=.48,below=2pt] {可能引起响动} (bush.east);
\draw[edge] (wind.east) to[out=5,in=-90] node[relation,pos=.66,below] {威胁很小} (self.south);
\draw[edge,densely dotted] (self.east) to[out=12,in=-92] node[relation,pos=.55,right] {身处} (america.south);
\draw[edge,densely dotted] (america.north west) to[out=155,in=90,looseness=1.14]
  node[relation,pos=.52,above] {在其自然分布区以外} (tiger.north);
\end{tikzpicture}

\caption{用知识图谱（knowledge graph）表示对一个人（自身）可能具有意义的信息，例子是“自身”察觉一丛灌木沙沙作响，因而要判断可能受到什么威胁。实线边和实线边框表示自身已有的知识，点线边和点线边框表示自身可能学到的新知识或新含义。学到的信息可帮助自身就所处环境中的个人风险作出很有用的推断。语言必须能够精确编码这样的复杂关系，才能尽可能发挥实际效用，同时尽量减少可能致命的误解。}
\label{fig:ch01-graph}
\end{figure}

把一生中经历的事抽象成带有意义的概念，是人类与世界发生有利于自身演化的互动的基础。以灌木沙沙作响为典型例子。原因未知时，图~\ref{fig:ch01-graph} 展示了与这一现象相关的部分知识可能如何编码在一个人的大脑里。“自身”附近的灌木发出响声，可能有多种含义。如果经过一连串推理，认为响声并非源于某种生存威胁，就可以判断这一观察与生存无关。另一方面，如果只根据图~\ref{fig:ch01-graph} 知识图谱中的实体及其关系，而风又不足以造成这么大的响动，那么认定响动来自老虎可能更安全。

能够如此推断外部世界，显然并不以人类的自然语言为必要条件，许多其他动物似乎也有类似的高级认知能力。语言使具有实际作用的意义单元能够传给任何能解码其中信息的个体。于是，有助于生存的知识和意义可以在人之间广泛传播，摆脱亲身经历的限制，也突破了低容量、低精度交流方式的限制。这在动物界前所未有，是人类取得生态优势的基础。

为了完成这项难事，人类的语言系统演化得十分复杂。语言既要编码关于世界的明确事实，也要能把事实嵌入不同的概念情境。强大的大脑在处理事实时，不同情境往往会改变事实的含义或重要性。语言若做不到这一点，就难以足够精确地传达信息，让接收者将它用于有价值的行动。直到不久前，只有人类拥有创造和使用人类语言所需的精密硬件，也可以称作生物湿件（wetware）。在人工智能领域，机器学习与自然语言处理推动了理解和生成人类语言的变革，并催生了大语言模型。下面先回顾这段历史，再细看大语言模型的能力。

% PDF 36--41。
\section{自然语言处理与语言模型的演进}
在计算机科学领域，自然语言处理是一门重要学科，旨在促成机器与人类语言之间的交互。它研究如何设计算法和技术，使计算机通过自然语言理解（Natural Language Understanding，NLU）理解和解释人类语言，并通过自然语言生成（Natural Language Generation，NLG）生成近似人类表达的文本。人类语言无疑极其复杂，人们对语言理解算法的需求日益增加，也凸显了自然语言处理的重要性。过去，自然语言处理策略通常以语言学为中心，分析从词性等基本语义和句法成分入手。当代深度学习（deep learning）方法则可能省去这类中间成分，自行形成能用于广泛任务的表示。图~\ref{fig:ch01-timeline} 展示了自然语言处理，特别是语言模型发展过程中的开创性研究。本节将作简要回顾。

\subsection{基于句法和文法的方法（20 世纪 60 至 80 年代）}
早在 20 世纪 40 年代，这一领域就有初步探索。1954 年的 IBM—乔治城大学实验（IBM-Georgetown experiment）展示了用机器将大约 60 个句子从俄语译为英语，是该领域的重要里程碑（\refcite{ch01-hutchins}{Hutchins et al., 1955}）。20 世纪 50 年代末，一系列开创性工作改变了语言理解研究的面貌。1957 年，诺姆·乔姆斯基（Noam Chomsky）出版《句法结构》（\textit{Syntactic Structures}），是语言学研究的一个转折点。该书凸显了句子的句法结构对理解语言的重要性（\refcite{ch01-chomsky}{Chomsky, 1957}）。同期，短语结构文法（phrase-structure grammar）的出现进一步影响了当时语言学与计算研究的视角。人工智能也有一系列进展，例如约翰·麦卡锡（John McCarthy）于 1958 年创立 LISP，以及第一个获得广泛认可的聊天机器人 ELIZA 的开发。这些成果塑造了自然语言处理的演进，也深刻影响了更广泛的人工智能领域。SHRDLU 是一个早期系统，通过结合句法、语义与推理来辨识基本查询和应答。该时代的系统，如 ELIZA（1966 年）与 SHRDLU（1970 年），主要依靠预先规定的词表与规则集理解和生成语言。
\transnote{原书第 6 页将“短语结构文法”的英文印为 \textit{phase-structure grammar}。}

\begin{figure}[!htbp]
\centering
% 重绘自原书第 6 页的图 1.2。
\begin{tikzpicture}[
  font=\footnotesize\sffamily,
  band/.style={minimum height=6.5mm,minimum width=22.8mm,inner sep=0pt,text=white,font=\scriptsize\sffamily\bfseries},
  event/.style={text width=19.8mm,inner sep=1.3mm,rounded corners=1pt,
    align=left,fill=black!5,font=\scriptsize\sffamily},
  phase/.style={font=\footnotesize\sffamily\bfseries,text width=21mm,align=center,inner sep=1pt}]
\node[band,fill=red!55!black] (p1) at (0,0) {\hspace{2mm}1950--80 年代};
\node[band,fill=blue!48!black] (p2) at (2.28,0) {\hspace{2mm}1980--2000 年代};
\node[band,fill=orange!80!black] (p3) at (4.56,0) {\hspace{2mm}2000--10 年代};
\node[band,fill=violet!70!black] (p4) at (6.84,0) {\hspace{2mm}2010--20 年代};
\node[band,fill=yellow!62!orange,text=black] (p5) at (9.12,0) {\hspace{2mm}2020 年代起};

\foreach \period in {1,...,5} {
  \node[circle,fill=white,text=black!84,inner sep=.7pt,font=\tiny\sffamily\bfseries]
    at ([xshift=2.15mm]p\period.west) {\period};
}
\node[phase,below=2mm of p1] (t1) {基于句法\\和文法};
\node[event,below=2mm of t1] {1. 乔姆斯基的《句法结构》\\2. ELIZA 聊天机器人\\3. 基于规则的 SHRDLU};
\node[event,above=2mm of p2] (e2) {1. 基于规则与本体的系统\\2. 统计模型\\3. $n$ 元模型结合机器学习算法};
\node[phase,above=2mm of e2] {专家系统\\与统计模型};
\node[phase,below=2mm of p3] (t3) {神经模型\\与稠密表示};
\node[event,below=2mm of t3] {1. Bengio 等人的稠密向量表示\\2. Mikolov 等人的循环神经网络语言模型\\3. 预训练词嵌入};
\node[event,above=2mm of p4] (e4) {1. word2vec、GloVe 等词嵌入\\2. 预训练与微调式迁移学习\\3. Bahdanau 等人的注意力机制\\4. Vaswani 等人的 Transformer\\5. BERT、GPT 等模型};
\node[phase,above=2mm of e4] {深度学习革命};
\node[phase,below=2mm of p5] (t5) {大语言模型时代};
\node[event,below=2mm of t5] {1. OpenAI 发布 GPT-2、GPT-3.5 和 GPT-4\\2. 用 RLHF 与人类价值观对齐，如安全性和有据可依\\3. 开源大语言模型与框架};
\end{tikzpicture}

\caption{自然语言处理和大语言模型从 20 世纪 50 年代到本书写作时的演进时间线，标出主要事件和突破。}
\label{fig:ch01-timeline}
\end{figure}

\subsection{专家系统与统计模型（20 世纪 80 年代至 21 世纪初）}
20 世纪 80 年代初，自然语言处理主要采用符号方法，这类系统通常称为专家系统（expert system）。它们的特点是规则和本体（ontology）由人工建立，二者构成结构化知识库，详列特定领域内的事实、概念及相互联系。到 20 世纪 80 年代末，符号人工智能的局限已经显现，特别是难以扩展，也难以处理自然语言中的歧义。认识到这一点后，研究者逐步转向统计方法。其思想很简单，从数据学习模式，省去逐条写死所有可能规则的工作。宾州树库（Penn TreeBank）、英国国家语料库（British National Corpus）、布拉格依存树库（Prague Dependency Treebank）和 WordNet 等语言资源，包括一系列带标注的汇编资料，对学术研究和商业应用都极有价值（\refcite{ch01-hajicova}{Hajicová et al., 1999}；\refcite{ch01-marcus}{Marcus et al., 1994}；\refcite{ch01-miller}{Miller, 1995}）。20 世纪 80 年代引入的隐马尔可夫模型（Hidden Markov Model，HMM），是统计方法在语言上的最早一批成功应用之一。它将语言视为一系列状态及其转移，在早期自然语言处理中得到广泛应用，尤其是在语音识别系统中。

$n$ 元模型（$n$-gram model）成为统计语言建模的基础，为刻画语言的序列性质提供了一种基本方法。它以条件概率为基础，核心思想是根据前面 $n-1$ 个词构成的历史估计当前词出现的概率，可写为
\[
P(w_n\mid w_{n-1},w_{n-2},\ldots,w_1).
\]
但是，当 $n$ 很大时，直接估计这个概率的计算开销可能很大，也容易遇到数据稀疏问题。因此，$n$ 元模型往往采用马尔可夫假设（Markov assumption），将历史缩减为最近 $n-1$ 个词。
\[
P(w_n\mid w_{n-1},w_{n-2},\ldots,w_1)
\approx P(w_n\mid w_{n-1},w_{n-2},\ldots,w_{n-(n-1)}).
\]
\transnote{保留原书第 7 页的下标。原式同时用 $n$ 表示当前位置和窗口阶数，因而 $n-(n-1)=1$，式面上未截短历史。若当前位置写作 $t$，则 $n$ 元模型保留的最早位置是 $t-(n-1)$。}

对于一元模型（unigram model，$n=1$），不考虑前序词，独立估计每个词的概率，写作
\[
P(w_n).
\]
在这种情况下，假定每个词独立生成。相应概率等于这个词在语料库中的出现次数除以词的总数。对于二元模型（bigram model，$n=2$），一个词的概率以紧邻的前一个词为条件。
\[
P(w_n\mid w_{n-1}).
\]
对于三元模型（trigram model，$n=3$），一个词的概率以前面紧邻的两个词为条件。
\[
P(w_n\mid w_{n-1},w_{n-2}).
\]
图~\ref{fig:ch01-bigram} 演示了这种概率推断过程。

\begin{figure}[!htbp]
\centering
% 重绘自原书第 7 页的图 1.3。模型输入及候选词保留英文。
\begin{tikzpicture}[font=\small\sffamily,
  candidate/.style={draw=black!55,rounded corners=1pt,text width=49mm,align=left,minimum height=4.4mm,inner sep=1.3pt,fill=yellow!23}]
\node[text width=43mm,align=left,text=violet] (prompt) at (0,0) {ChatGPT is one of the\\large language \ldots};
\node[candidate,fill=green!14] (title) at (5.5,1.65) {概率（\textcolor{red!80!black}{词}\,$\mid$\,前一个词）};
\node[candidate] (models) at (5.5,.96) {$P(\text{\textcolor{red!80!black}{models}}\mid\text{language})=0.43$};
\node[candidate] at (5.5,.49) {$P(\text{\textcolor{red!80!black}{processing}}\mid\text{language})=0.21$};
\node[candidate] at (5.5,.02) {$P(\text{\textcolor{red!80!black}{translation}}\mid\text{language})=0.12$};
\node[font=\footnotesize] at (5.5,-.42) {\ldots\hspace{14mm}\ldots\hspace{14mm}\ldots};
\node[font=\footnotesize] at (5.5,-.76) {\ldots\hspace{14mm}\ldots\hspace{14mm}\ldots};
\node[candidate] (barrier) at (5.5,-1.2) {$P(\text{\textcolor{red!80!black}{barrier}}\mid\text{language})=0.04$};
\coordinate (split) at (2.52,0);
\draw[thick,violet] (prompt.east) -- (split);
\draw[thick,violet,-{Stealth[length=1.9mm]}] (split) |- (models.west);
\draw[thick,violet,-{Stealth[length=1.9mm]}] (split) |- (barrier.west);
\end{tikzpicture}

\caption{用假想的二元模型补全片段“ChatGPT is one of the large language”的过程。每个黄色方框给出一个可能接在“language”后的词及其条件概率。二元概率由训练语料计算得到。}
\label{fig:ch01-bigram}
\end{figure}

\begin{keypoint}
给定一段词序列，模型依据从自然语言句子组成的训练数据中学得的参数，预测后面各候选词出现的概率，并按概率排列最可能的续写。我们把采用这一概率框架的模型称为语言模型（Language Model，LM）或预训练语言模型（Pre-trained Language Model，PLM）。
\end{keypoint}

传统上，大多数自然语言处理工作将 $n$ 元方法和机器学习技术结合，在监督和无监督设置下处理大量自然语言处理任务。相关技术包括多项逻辑回归（multinomial logistic regression）、支持向量机（support vector machine）、贝叶斯网络（Bayesian network）和期望最大化算法（expectation-maximization algorithm）（\refcite{ch01-brown-1992}{Brown et al., 1992}；\refcite{ch01-manning}{Manning and Schütze, 1999}）。

\subsection{神经模型与稠密表示（21 世纪头两个十年）}
21 世纪初，\refcite{ch01-bengio}{Bengio 等人（2000）}的开创性研究产生了第一个神经语言模型（neural language model）。该模型用查找表（lookup table）映射前面的 $n$ 个词，再将它们送入带隐藏层的前馈网络（feed-forward network），最后由 softmax 层对输出作平滑处理，以预测下一个词。这项研究的重要意义在于，它改变了传统的 $n$ 元模型或词袋模型（bag-of-words model）的做法，把“稠密向量表示”（dense vector representation）引入自然语言处理。此后，采用循环神经网络（Recurrent Neural Network，RNN）和长短期记忆（Long Short-Term Memory，LSTM）架构的语言模型，成为该领域的前沿方案（\refcite{ch01-graves}{Graves, 2013}；\refcite{ch01-mikolov-2010}{Mikolov et al., 2010}）。\refcite{ch01-collobert}{Collobert 和 Weston（2008）}的重要研究初步展示了预训练词嵌入（pre-trained word embedding）以及用卷积神经网络（Convolutional Neural Network，CNN）分析文本的做法。他们还强调，查找表，也就是如今所称的嵌入矩阵（embedding matrix），在多任务学习（multitask learning）中很有用。

\subsection{深度学习革命（2010 年代至 2020 年代）}
词嵌入技术的重要进步得益于 \refcite{ch01-mikolov-2013a}{Mikolov 等人（2013a}、\refcite{ch01-mikolov-2013b}{b）}的工作，他们提高了 \refcite{ch01-bengio}{Bengio 等人（2000）}最初构想的嵌入的训练效率。他们去掉隐藏层并采用近似学习目标，提出了可高效、大规模训练词嵌入的 word2vec。它有两个变体。连续词袋（Continuous Bag-of-Words，CBOW）根据周围的词预测中心词，跳元模型（skip-gram）反过来预测邻近词。训练效率的提升使在大规模语料库上训练成为可能，使稠密嵌入能够涵盖多种语义关系。以这类词嵌入作为初始表示，再在大规模数据集上预训练，已经成为神经网络架构的一项基本方法。

图~\ref{fig:ch01-word2vec} 直观演示了 word2vec 的过程，重点展示如何把词变成高维嵌入。单个词先被映射到高维空间中相应的嵌入，再在低维的二维平面上显示。

\sourcefigure{fig01_04.pdf}{这个示例展示词的表示与词在语义空间中相似性的关系。从“man”到“woman”的向量与从“king”到“queen”的向量十分相似。这说明在嵌入空间中，“man”与“woman”之间的关系（或差别），类似于“king”与“queen”之间的关系（或差别）。另外，“dog”靠近“cat”和“kitten”，而“houses”与这些动物词明显分开。原图从左到右为词、词嵌入、从七维到二维的降维和二维可视化。}{fig:ch01-word2vec}

自然语言处理和语音处理所用神经架构的另一个重要进展，是 \refcite{ch01-bahdanau}{Bahdanau 等人（2014）}提出的注意力机制。随后，\refcite{ch01-vaswani}{Vaswani 等人（2017）}于 2017 年提出 Transformer 架构，改变了语言模型的基本范式。Transformer 利用注意力机制，按输入数据的各部分与任务的相关程度，有选择地加以关注，因此在多项自然语言处理任务上显著提高了表现。这个基础架构为 BERT、GPT-1 等后继模型铺平了道路。Google 于 2018 年推出 BERT，标志着迁移学习（transfer learning）在自然语言处理领域取得重要进展。BERT 先在大规模文本语料上训练，再针对特定任务微调。先让模型获得一般能力，再针对任务微调，成为自然语言处理中的常见做法。

% PDF 41--44。
\section{大语言模型时代}
深度学习革命之后，基于 Transformer 架构的小型语言模型取得显著成果，大语言模型出现的条件也已具备。这类新型语言模型的潜力十分明显，它改变了机器学习从业者解决常见自然语言处理问题的方式。为了利用大语言模型前所未有的语言理解与文本生成能力，世界各地的开发和研究团队正将文本分类（text classification）、命名实体识别（Named Entity Recognition，NER）等长期存在的语言任务重新表述为文本生成任务（\refcite{ch01-zhao}{Zhao et al., 2023}）。

\begin{practicaltip}
从模型架构看，大语言模型与基于 Transformer 的较小语言模型或预训练语言模型，主要以参数数量区别。两者之间没有公认的界限，但大语言模型通常有数亿到数万亿个参数。GPT-1 与 BERT 等最早的 Transformer 模型可以视作第一代大语言模型。特意设计得较小的模型，通常用于算力有限的场景，可视作小型预训练语言模型。
\end{practicaltip}

\subsection{大语言模型演进简史}
2022 年 11 月发布 ChatGPT，无疑是大语言模型研究的一个转折点。大语言模型的能力首次如此广泛地引起公众关注。然而，要解释我们如何走到今天，不能只看这一个事件。下面回顾在 ChatGPT 发布之前和之后推动大语言模型发展的其他重要事件。

注意力机制与 Transformer 架构问世之后，T5 的发布是一次重要创新，让人们初步看到了当代大语言模型的潜力（\refcite{ch01-raffel}{Raffel et al., 2020}）。T5 是一个文本到文本的语言模型，完全不使用循环或卷积网络结构，只利用 Transformer 的注意力网络。得益于模型规模（当时有 110 亿个参数）和预训练语料规模（1 万亿个词元，token），T5 在多项文本到文本的翻译任务中达到了当时的先进水平。T5 发布时还给出了多种自然语言处理任务的评估结果。这些任务都重新表述为文本到文本的问题，模型表现出色。Google 发布 T5 后不久，OpenAI 也推出 GPT-3，作为此前 GPT-1 和 GPT-2 语言模型的后继（\refcite{ch01-brown-2020}{Brown et al., 2020}）。GPT-3 显示出很强的泛化能力，包括出色的上下文学习能力（见第~\ref{sec:ch01-icl} 节）。这些能力来自无监督预训练，无需针对特定任务再微调。

从 2020 年初夏 GPT-3 发布到 2022 年 11 月 ChatGPT 发布，其间发生了几件重要的事。一个重要里程碑是 Anthropic 发布了拥有 520 亿个参数的模型，首次提出使大语言模型的回答更符合人类价值观的方法（\refcite{ch01-askell}{Askell et al., 2021}）。在此之前，GPT-3 等模型的用户已经清楚看到一个根本局限，即模型回答往往不符合人类层面的价值观。这项工作显示，对齐调优有望让模型在对话中表现得更像人类。不久后，OpenAI 推出另一款模型 InstructGPT（\refcite{ch01-ouyang}{Ouyang et al., 2022}）。这是 GPT-3 系列的一个微调变体，在对齐阶段使用人类反馈进行调优。对齐之后，InstructGPT 更善于给出人类评估者偏好的回答。即使只有 13 亿个参数，其表现也优于拥有 1\,750 亿个参数的 GPT-3。这表明，尽管大语言模型要表现出某些最有用的涌现能力，必须扩大参数规模，但微调也是使语言模型在多项任务上取得更好表现的一条途径。这为 ChatGPT 的诞生奠定了基础。

自 2022 年 11 月 ChatGPT 发布以来，OpenAI 又发布了 GPT-4（\refcite{ch01-openai}{OpenAI, 2023}）。在标准的自然语言处理和自然语言理解基准上，GPT-4 的能力比最初版 ChatGPT 更强。它在律师资格考试等面向人类的测试中表现明显改善（\refcite{ch01-martinez}{Martínez, 2023}），也获得了理解图像与文本之间关系的新能力。除 OpenAI 的成果外，写作时其他值得关注的大语言模型包括 Meta 的 Llama-2（\refcite{ch01-touvron}{Touvron et al., 2023}）、Google Research 团队的 PaLM\footnote{\url{https://blog.google/technology/ai/google-palm-2-ai-large-language-model}} 和 Anthropic 的 Claude-2。\footnote{\url{https://www.anthropic.com/index/claude-2}}

读者完全可以预期，大语言模型研究的创新及其应用还将继续发展。这些技术带来前所未有的人机交互机会，也是通过计算把人类意图和目标大规模付诸实践的最有前景的途径之一。无论是更高效、更全面地解决传统自然语言处理问题，还是开拓以往不可能的应用，作者都期待与读者一道深入这一领域。

\subsection{大语言模型的规模}
任何大语言模型的规模都涉及三个主要因素。
\begin{enumerate}
\item 预训练语料规模，决定训练时模型习得知识的广度和深度。
\item 学得的参数数量，决定模型学得的状态有多复杂。
\item 计算规模，关系到采用给定架构进行训练和推理是否可行。
\end{enumerate}
已有大量研究试图理解，扩大这三个因素的规模如何带来所谓的涌现能力（emergent ability）。所谓涌现，是指大语言模型出现了在预训练及后续微调中都没有明确训练过的能力（\refcite{ch01-hoffmann}{Hoffmann et al., 2022}；\refcite{ch01-wei-2022}{Wei et al., 2022}）。

\subsection{大语言模型的涌现能力}
讨论涌现能力之前，可以先看看何种情形不算涌现。早期不断扩大模型规模的工作，确实在许多自然语言处理问题上取得了成果。例如，BERT 在实体关系分类（entity-relation classification）任务中表现突出（\refcite{ch01-soares}{Soares et al., 2019}）。在这类应用中，较小的预训练语言模型也有很好的表现。不过，提升通常来自微调过程。研究者精心设计目标，投入大量精力整理能体现这个目标的数据集，例如带标签的样例或问答对，再对预训练语言模型的参数空间进一步调优。在这种情形下，模型获得的有用能力是明确教给它的。

涌现能力则无需这些额外的微调步骤，甚至不必明确教给模型。具有涌现能力的大语言模型可以在完全不改变预训练权重的情况下“学会”解决这类问题（\refcite{ch01-wei-2022}{Wei et al., 2022}）。如果模型够大、训练语料覆盖够广、训练算力预算够充足，大语言模型就会在具体自然语言处理任务和语言理解、算术、多步推理等更高层次的任务上开始表现出很强的能力（\refcite{ch01-radford}{Radford et al., 2019}；\refcite{ch01-wei-2022}{Wei et al., 2022}）。

\sourcefigure{fig01_05.pdf}{不同规模的大语言模型表现出的涌现能力。横轴以模型参数数量表示规模，纵轴给出任务准确率（B 为精确匹配率）。在少样本设置下，评估 LaMDA（\refcite{ch01-thoppilan}{Thoppilan et al., 2022}）、GPT-3（\refcite{ch01-brown-2020}{Brown et al., 2020}）、Gopher（\refcite{ch01-rae}{Rae et al., 2022}）、Chinchilla（\refcite{ch01-hoffmann}{Hoffmann et al., 2022}）和 PaLM（\refcite{ch01-chowdhery}{Chowdhery et al., 2022}）在四项不同的自然语言处理任务上的表现。A 和 B 来自 BIG-Bench 基准套件（\refcite{ch01-srivastava}{Srivastava et al., 2023}）。A 测试三位数加减和两位数乘法，B 测试波斯语问答，C 综合 57 项不同任务的表现，D 测试语义理解。每项任务中，至少有一个大语言模型表现出明显的能力涌现趋势。}{fig:ch01-emergence}

图~\ref{fig:ch01-emergence} 展示了扩大大语言模型规模对四项自然语言处理任务准确率的影响。在每一项任务中，较小的语言模型都不优于随机作答，但超过一定参数规模后，至少有一个模型开始明显胜过随机作答。图~\ref{fig:ch01-emergence} 部分重制自 \refcite{ch01-wei-2022}{Wei 等人（2022）}。图中准确率往往突然上升，表明仅凭参数较少的大语言模型的表现，无法预先判断这些能力会涌现。涌现能力怎样或为何出现，仍是活跃的研究方向。

% PDF 44--55；原书第 13--24 页。
\section{大语言模型实践}
像任何新技术一样，大语言模型在走向成熟乃至普及的过程中，研究与应用阶段有很大的重叠。因而，找到最有价值或最有用的研究文献，或者为某个应用选出最合适的方法，都可能颇为困难。下面几节将梳理大语言模型在实际工作中的不同用法，帮助读者作出判断。

\begin{figure}[!htbp]
\centering
% 重绘自原书第 14 页的图 1.6。
\begin{tikzpicture}[
  font=\footnotesize\sffamily,
  flow/.style={draw=black!50,-{Stealth[length=1.6mm]},line width=.55pt},
  choice/.style={diamond,draw=black!40,aspect=1.54,minimum height=12mm,align=center,inner sep=1.4pt},
  process/.style={draw=black!40,rounded corners=2mm,fill=green!22!gray!18,
    minimum width=18mm,minimum height=9mm,text width=16mm,align=center,inner sep=1pt},
  stage/.style={font=\footnotesize\sffamily\bfseries},
  app/.style={align=center,inner sep=1pt}]
\node[stage] (scope) at (.56,5.75) {确定范围};
\node[stage] (pretrain) at (3.0,5.75) {预训练};
\node[stage] (adapt) at (6.53,5.75) {适配与对齐};
\node[stage] (deploy) at (9.95,5.75) {应用};
\draw[flow,dashed] (scope.east) -- (pretrain.west);
\draw[flow,dashed] (pretrain.east) -- (adapt.west);
\draw[flow,dashed] (adapt.east) -- (deploy.west);

\node[choice,fill=orange!50] (llm) at (.65,2.74) {模型？};
\node[process,fill=yellow!25,text width=19mm] (scratch) at (2.65,4.55)
  {从头预训练\\[-1pt]{\scriptsize 图书、网页语料、代码}};
\node[process,fill=yellow!45,text width=19mm] (ready) at (2.65,.76)
  {预训练大语言模型\\[-1pt]{\scriptsize OpenAI、Llama、Claude}};
\node[choice,fill=green!18!gray!14] (need) at (4.34,2.74) {适配？};
\node[process] (prompt) at (6.47,4.55) {提示工程};
\node[process] (finetune) at (6.47,2.74) {微调};
\node[process] (align) at (6.47,.87) {对齐调优};
\node[choice,fill=blue!16!gray!10,aspect=1.66] (augment) at (8.85,2.74) {增强？};
\node[app] (user) at (10.55,4.55) {查询提示};
\node[process,fill=blue!37!cyan!30,minimum width=9mm,text width=8mm] (rag) at (10.55,3.45) {RAG};
\node[app,draw=blue!55!black,cylinder,shape border rotate=90,
  aspect=.35,minimum height=7mm,minimum width=12mm,fill=blue!12] (db) at (10.55,2.06) {向量库};
\node[app,draw=blue!55!black,minimum height=7mm] (docs) at (10.55,.7) {自定义文档};

\draw[flow] (llm.north) |- (scratch.west);
\draw[flow] (llm.south) |- (ready.west);
\draw[flow] (scratch.east) -| (need.north);
\draw[flow] (ready.east) -| (need.south);
\coordinate (split) at (5.40,2.74);
\draw[draw=black!50,line width=.55pt] (need.east) -- (split);
\draw[flow] (split) |- (prompt.west);
\draw[flow] (split) -- (finetune.west);
\draw[flow] (split) |- (align.west);
\draw[flow] (prompt.south) -- (finetune.north);
\draw[flow] (finetune.south) -- (align.north);
\draw[flow] (prompt.east) -| (augment.north);
\draw[flow] (finetune.east) -- (augment.west);
\draw[flow] (align.east) -| (augment.south);
\draw[flow] (augment.east) -- (10.55,2.74) -- (rag.south);
\draw[flow] (user.south) -- (rag.north);
\draw[flow] (db.north) -- (10.55,2.74);
\draw[flow] (docs.north) -- (db.south);
\end{tikzpicture}

\caption{从应用出发规划大语言模型的开发、适配和使用时，可以考虑的各种路径。}
\label{fig:ch01-paths}
\end{figure}

大语言模型的实践大体涉及三个领域。
\begin{itemize}
\item \textbf{开发（development）}关注如何从头构建大语言模型，包括从数据集、学习目标和网络架构几个角度考虑预训练。
\item \textbf{适配（adaptation）}关注如何修改预训练的大语言模型，使其产生更符合预期的结果。这些结果往往要结合具体应用衡量，应用也会影响选择哪种适配方法。指令调优、对齐和微调是这里的重要主题。
\item \textbf{使用（utilization）}关注如何与大语言模型交互，以发挥其有用的能力。终端用户提示（end-user prompting）和应用开发是这里的重点。
\end{itemize}

本节先概述大语言模型开发，第 2 章再作详细讨论。接着介绍最常用的适配方法，包括对齐调优和参数高效微调，第 4 和第 5 章会展开这些主题。最后，从使用大语言模型的角度考察终端用户提示和应用。后文还会详谈这些内容，相关材料主要见第 6、7、8 章。下面分别介绍每一领域的基本概念。

\subsection{大语言模型开发}
可以从头预训练自己的大语言模型，也可以在已有模型上进行适应性预训练（adaptive pre-training）。由于这些做法复杂，前期成本也高，预计它们不会成为特别常见的开发路径。若团队需要模型掌握专业领域知识，希望更明确地控制偏见、能力或局限，或者希望降低长期的总拥有成本（Total Cost of Ownership，TCO），就可能选择这条路径，而不用已有的众多开源或闭源预训练模型。本节介绍这些做法所涉及的主要因素。

\subsubsection{大语言模型预训练}
大语言模型的许多出色能力来自预训练过程。预训练使模型从大规模数据集中习得基本的语言技能和知识（\refcite{ch01-zhao}{Zhao et al., 2023}）。下面回顾预训练的三个方面，即常用的数据集、网络架构和学习目标。
\begin{itemize}
\item \textbf{数据集。}与开发大语言模型所用的学习目标和网络架构相比，数据集在不同模型间的差异较小。训练通用大语言模型时，通常组合使用网页抓取资料（如 Common Crawl）、维基百科（Wikipedia）、图书语料和代码。代码数据已被证明能显著提高模型的推理能力（\refcite{ch01-ma}{Ma et al., 2023}）。这些来源不同的内容还要经过筛选整理，统一质量标准，并尽量减轻数据重复、有害或仇恨内容等因素的负面影响。确认数据集覆盖范围和质量足够后，开发团队按需要对数据作规范化处理和词元切分（tokenization）。完成这些步骤，数据才可用于预训练。第 2.4 节会更详细地讨论这些内容。

\item \textbf{网络架构。}大语言模型最常用的是 Transformer 架构，它在开创性的论文《注意力就是你所需要的一切》（\textit{Attention Is All You Need}）中提出（\refcite{ch01-vaswani}{Vaswani et al., 2017}）。原始架构由编码器和解码器组成，两者都采用自注意力机制（self-attention mechanism）和前馈神经网络（feed-forward neural network）。第 2.5.3 节将详细介绍研究者怎样改进原始 Transformer 架构。这里选取几个影响较大的模型，说明架构的改进与模型能力之间的关系。
\begin{itemize}[label=--]
\item 纯编码器架构（encoder-only architecture）用于得到输入序列的上下文表示（contextual representation）。这种表示对分类、命名实体识别等多种自然语言处理任务很有用。BERT 是采用这种设计的一个常用模型（\refcite{ch01-devlin}{Devlin et al., 2019}）。
\item 编码器—解码器架构用于机器翻译、问答等序列到序列（sequence-to-sequence）任务。T5 是采用这种架构的一个常用模型（\refcite{ch01-raffel}{Raffel et al., 2020}）。
\item 当学习目标是自回归序列生成时，使用因果解码器架构（causal decoder architecture）。将注意力限制在一个方向上，就可以实现序列生成。采用这种架构的模型擅长文本生成，最熟悉的是 GPT 系列（例如 \refcite{ch01-brown-2020}{Brown et al., 2020}）。
\item 前缀解码器架构（prefix-decoder architecture），又称非因果解码器架构（non-causal decoder architecture），是上述因果解码器的变体。关键差异在于，对输入序列，即前缀，应用双向注意力机制。生成的序列仍采用单向注意力，生成过程仍是自回归的。书中以 Google 的 PaLM 为例，说明这种架构尤其擅长机器翻译等可受益于双向编码的任务（\refcite{ch01-chowdhery}{Chowdhery et al., 2022}）。
\end{itemize}

\item \textbf{学习目标。}预训练所用的学习目标，对大语言模型最终表现出的涌现能力有很大影响。通常，预训练目标旨在尽可能提高自然语言理解和连贯生成的能力。最常用的方法是完整语言建模（full language modeling），即给定此前的词元，以自回归方式预测序列中的下一个词元（\refcite{ch01-zhao}{Zhao et al., 2023}）。其他重要目标包括去噪自编码（denoising autoencoding）与混合去噪器（mixture-of-denoisers）。前者先破坏输入文本序列，再训练网络恢复受破坏的片段（\refcite{ch01-raffel}{Raffel et al., 2020}）。后者组合使用擅长不同任务的去噪器，利用三种主要的预训练任务，分别为标准语言建模、恢复带少量噪声的短片段，以及恢复长片段或高噪声片段（\refcite{ch01-chowdhery}{Chowdhery et al., 2022}）。前者用于常见的序列到序列大语言模型 T5，原书将后者列为 Google 的 PaLM 所用方法。
\end{itemize}

\subsubsection{大语言模型的适应性预训练}
大语言模型通常使用上一节所述的通用数据集预训练，而领域适应性预训练（domain-adaptive pre-training）使用更贴近特定领域的数据集继续训练模型（\refcite{ch01-gururangan}{Gururangan et al., 2020}）。它的目标是使大语言模型的能力更适应有专门词汇或语言用法的领域。例如，生物学有大量术语，物种名称和解剖描述还会使用拉丁语命名法。只在通用数据集上训练的大语言模型，可能对生物学术语掌握得不够充分。使用来自 bioRxiv 之类的数据集适配基础模型，可以克服这些局限。bioRxiv 是面向生物科学文献的专业预印本平台（\refcite{ch01-zhang-a}{Zhang et al., 2024a}）。

\begin{figure}[!htbp]
\centering
% 重绘自原书第 17 页的图 1.7。
\begin{tikzpicture}[font=\small\sffamily]
\begin{scope}
\clip (-3,-1.75) rectangle (3,2.35);
\fill[cyan!30] (-2.23,.76) ellipse [x radius=2.94cm,y radius=1.86cm,rotate=-28];
\fill[yellow!72!orange,opacity=.36] (2.42,-.20) ellipse [x radius=2.88cm,y radius=1.68cm,rotate=3];
\fill[gray!55,opacity=.62] (-.6,.67) ellipse [x radius=1.30cm,y radius=.60cm,rotate=-50];
\fill[black] (-.54,.41) ellipse [x radius=.65cm,y radius=.22cm,rotate=-39];
\end{scope}
\node at (-1.76,-.74) {目标领域};
\node[align=center] at (1.78,-.65) {原大语言模型\\领域};
\node (task) at (.14,2.00) {任务};
\draw[-{Stealth[length=1.9mm]},black!70] (task.south) to[out=-92,in=84] (-.59,.61);
\end{tikzpicture}

\caption{通用大语言模型预训练（原模型领域）、领域适应性预训练（目标领域）和任务微调（任务）涉及的数据分布之间的关系。任务分布所处的浅灰色区域表明，任务通常是较大分布中的一个可观测子集。}
\label{fig:ch01-domain}
\end{figure}

研究还表明，适应性预训练可以显著改善任务微调的效果（\refcite{ch01-gururangan}{Gururangan et al., 2020}），例如特定领域的主题分类，见图~\ref{fig:ch01-domain}。多个领域都观察到了这种效果，金融服务业的一个知名例子是 BloombergGPT。该模型使用混合数据作适应性预训练，其中有新闻报道、财务报告等金融领域数据，也有 Bloomberg 的专有数据。尽管这部分数据只占所用预训练数据总量的 0.7\%，BloombergGPT 在金融领域的多项任务上仍明显优于参数量为它三倍的其他通用大语言模型（\refcite{ch01-wu}{Wu et al., 2023}）。
\transnote{原书第 17 页此处给出 0.7\%，未明确说明分子对应哪一部分数据，也未定义所比较的预训练数据总量。}

具备合适的数据集、足够的算力和充足的预算时，通用预训练与领域适应性预训练各有收益。只要能承担前期成本，以这种方式从头构建大语言模型，就可能在控制输出结果、性能或隐私方面获得显著优势。多数开发团队恐怕很难满足这一条件。下一节将介绍更接近微调的适配方法。这些方法通常只需上述预训练方法所需数据和计算成本的一小部分。

\subsection{大语言模型适配}
大语言模型的涌现能力本身令人瞩目，但即使经过适应性预训练，模型在给定任务上的开箱表现也未必满足每一种应用需求。围绕如何进一步调优大语言模型、使其朝指定结果改进的工具、技术和流程，研究迅速为开发者提供了许多选择。有些方法针对分类或命名实体识别等传统自然语言处理任务进行微调，有些则着重调整生成文本的风格，使之符合开发者基于价值判断形成的审美偏好。无论哪种，和前述预训练方法的主要区别都是使用带标签的训练数据集。本节探讨监督式适配的四个主要方面，包括指令调优、对齐调优、全参数微调和参数高效微调。

\subsubsection{指令调优}
指令调优是一种微调技术，用来增强大语言模型的能力和可控性。它的主要目标是微调通用大语言模型，使之更准确地遵循用户给出的具体指令（\refcite{ch01-zhang-b}{Zhang et al., 2024b}）。做法是将带标签的数据集整理成“（指令，输出）”对，对大语言模型进行监督训练。代码生成、摘要、问答、任务规划或执行等都可以整理成指令调优数据，使开发者能够提高模型在这些任务中遵循指令的能力。

指令调优通常通过全参数微调或参数高效微调来实现。在学习过程中，大语言模型更好地建立起指令与数据对中预期输出之间的联系，生成的回答因此更符合人类指令。不难预料，全参数指令调优的成本可能很高，特别是在用大型指令调优数据集训练超大模型时。因此，LoRA（\refcite{ch01-hu}{Hu et al., 2021}）等参数高效微调方法，以及专门用于指令调优的 HINT（\refcite{ch01-ivison}{Ivison et al., 2023}），都成为可行的方法，可以在所用数据集规模与微调成本之间取得更好的折衷。

\subsubsection{对齐调优}
指令调优旨在让大语言模型遵循人类的具体指令，而对齐调优旨在更广泛地使模型输出符合人类偏好和价值观（\refcite{ch01-lin}{Lin et al., 2023}）。和指令调优一样，原书将对齐调优归为监督微调技术。它高度依赖人工标注者，根据模型回答符合预先确定的偏好或价值观的程度，为回答排序。这涉及标注者本身的主观判断，尤其当他们来自不同的文化或社会背景时。要为微调过程提供一致的对齐反馈，这是一个重要难题，因此必须清楚界定我们希望维护的价值观。对齐调优通常采用以下三个原则。
\begin{itemize}
\item \textbf{有帮助性（helpfulness）}指模型紧扣提示中的指令，帮助用户完成任务的能力。
\item \textbf{诚实性（honesty）}指模型向用户提供准确信息的能力。
\item \textbf{无害性（harmlessness）}指模型避免生成会伤害用户或违背用户价值观与道德观的文本的能力。生成仇恨内容或表现出偏见等，都会降低模型的无害性。
\end{itemize}

第 5 章将详细讨论多种对齐调优方法。其中最知名的两种是基于人类反馈的强化学习（RLHF，见 \refcite{ch01-kaufmann}{Kaufmann et al., 2024}）和直接偏好优化（DPO，见 \refcite{ch01-rafailov}{Rafailov et al., 2023}）。RLHF 使用基于人类判断的反馈拟合反映人类偏好的奖励模型（reward model），再利用该模型微调大语言模型，使奖励最大化。拟合奖励模型和维持对齐训练的稳定性存在困难，DPO 的提出正是为了解决其中一些问题。它直接使用成对的偏好数据，按简单的分类目标（例如最大似然）微调大语言模型。提出者将 DPO 作为一种更简单的对齐调优方法。

\subsubsection{全参数微调}
早期的大语言模型微调主要采用指令调优。研究者将带标签的数据集重新整理成自然语言指令，让大语言模型处理这些指令，更新模型参数（见第 1.5.2.1 节）。然而，用数千条指令样本更新数十亿个参数，需要极高的算力，只有少数拥有足够预算的企业能承担。因此，人们投入很多精力，研究如何提高全参数调优的内存效率和计算效率。提高计算效率最常用的方法是量化（quantization），它在预训练期间或之后减少模型需要的内存（\refcite{ch01-gholami}{Gholami et al., 2021}）。除了用量化提高微调效率，近期研究也表明，降低优化过程的内存消耗既能帮助完成任务，也能减少微调的总成本（\refcite{ch01-lv}{Lv et al., 2023}）。第 4.4 节将进一步讨论这些方法。

\subsubsection{参数高效微调}
另一组大语言模型微调方法试图在提高特定任务表现的同时，尽量减少需要调优的参数数目（\refcite{ch01-zhao}{Zhao et al., 2023}）。下面介绍两种重要的参数高效微调方法。

\textbf{低秩适配（LoRA）}是 \refcite{ch01-hu}{Hu 等人（2021）}提出的方法。它在大语言模型的每个 Transformer 层中注入可训练的秩分解矩阵，同时冻结相应的预训练模型权重，从而减少需要调优的参数数目。注入的矩阵可以提高微调后任务的表现，又不会显著增加推理延迟。下面将看到，其他参数高效微调方法并非都能做到这一点。

\textbf{适配器}是另一种参数高效微调方法。它在每个 Transformer 层的输入和输出部件之间插入小型神经网络模块，以这些模块代替秩分解矩阵。优化适配器的参数时，规模大得多的 Transformer 部件保持固定。适配器显著减少需要调优的参数总量，因而大幅缩短训练时间。不过，在处理流程中加入额外部件会延长推理时间。第 4 章将更全面地讨论这些方法以及其他值得关注的参数高效微调方法。

\subsection{大语言模型的使用}
从最直接的层面看，使用大语言模型主要指终端用户通过提示与模型交互。不过，本节也从应用的层面讨论如何使用大语言模型，即从更高层面利用它。第 8 章会详细介绍这种用法。这里先介绍大语言模型中很有用的一种涌现能力，即上下文学习（\refcite{ch01-wei-2022}{Wei et al., 2022}），再介绍名为思维链提示的进阶方法，说明终端用户编写提示时的一些基本思路。提示工程还有许多创新，本节不逐一讨论，余下各章会介绍细节。从应用角度，本节概述几类主要应用，以及促成这些应用的概念和框架创新。本书后续各章将详谈其中许多主题。第 7 章讨论对话式大语言模型与检索增强生成，第 6 章讨论大语言模型的挑战与评估，第 8 章讨论大语言模型应用的开发与运维。

\subsubsection{上下文学习}
\label{sec:ch01-icl}
为了让大语言模型在任务中发挥涌现能力，往往会借助另一项涌现能力，即上下文学习。这种能力最早在 OpenAI 的早期 GPT 模型上得到展示（\refcite{ch01-brown-2020}{Brown et al., 2020}；\refcite{ch01-radford}{Radford et al., 2019}）。它使大语言模型仅在推理阶段就能从自然语言输入中学习，模型参数无须更新。通常把这类自然语言输入称为提示。根据输入提示中有多少任务示范作为上下文，可以分为零样本提示（zero-shot prompt）、少样本提示（few-shot prompt）或多样本提示（multi-shot prompt）。

\sourcefigure{fig01_08.pdf}{上下文学习演示。对三种参数数量不同的 GPT-3 模型（13 亿、130 亿和 1\,750 亿个参数）分别进行评估，测量它们移除故意插入单词的随机符号时的表现。每个模型都在零样本、单样本和少样本设置下测试。横轴为上下文中的示例数 $k$，纵轴为准确率。虚线表示提示中没有自然语言任务描述时的表现，实线表示把自然语言任务描述作为上下文时的表现。}{fig:ch01-icl}

图~\ref{fig:ch01-icl} 改编自 \refcite{ch01-brown-2020}{Brown 等人（2020）}的图 1.2。它展示了上下文学习对模型完成一项简单任务的影响，即去掉正常单词中插入的随机符号。在被测试的最大 GPT-3 模型（1\,750 亿个参数）上，这种影响最为显著。只给模型一条含自然语言任务描述的提示，就能明显改善表现。若在描述之外仅再给一个任务示范，相对于没有描述的情形，GPT-3 的表现会有更显著的跃升。GPT-3.5、GPT-4、PaLM、Llama-2 等先进模型的这一特性，是持续推动大语言模型应用创新的主要因素之一。怎样设计提示来改善语言模型在指定任务上的表现，第 3 章有详细讨论。原书还将读者引向与终端用户提示相关的第 7 章。
\transnote{原书第 21 页此处指向第 7 章。题为“终端用户提示”的专节是第 4.5 节。}

\subsubsection{思维链推理}
上下文学习的一个值得关注的特例是思维链提示（Chain-of-Thought Prompting，CoT）。这种方法组织提示中的上下文、指令和查询，引导大语言模型按有逻辑的一系列步骤来解决一个较大的任务（\refcite{ch01-wei-2023}{Wei et al., 2023}）。图~\ref{fig:ch01-cot} 演示了思维链的思想。与标准提示相比，在示范中加入逐步推理，可以让大语言模型更好地回答示例问题。上下文学习及其特例思维链提示，都属于在推理阶段用自然语言提示来引出大语言模型的知识、语言理解与推理能力的做法。

\begin{figure}[!htbp]
\centering
% 重绘自原书第 22 页的图 1.9。两个版本都保留例题与两个回答中的数值。
\begingroup
\small
\tcbset{cotpanel/.style={arc=2mm,colback=white,colframe=black!40,
  colbacktitle=black!82,coltitle=white,boxrule=.45pt,
  fonttitle=\sffamily\bfseries\small,left=4pt,right=4pt,top=4pt,bottom=4pt,
  before skip=5pt,after skip=5pt}}
\begin{minipage}[t]{.466\linewidth}
\centering{\sffamily\bfseries 标准提示}\par
\begin{tcolorbox}[cotpanel,title={模型输入}]
\textbf{问}\enspace Bob 举办聚会，买了 3 张披萨，每张切成 12 片。参加聚会的人吃掉了两整张披萨，以及另一张的三分之一。还剩多少片披萨？

\textbf{答}\enspace 答案是 8。

\textbf{问}\enspace Rebecca 举办聚会，买了 4 盒甜甜圈，每盒 12 个。聚会上有 12 人，每人吃 2 个甜甜圈。还剩多少盒甜甜圈？
\end{tcolorbox}
\begin{tcolorbox}[cotpanel,title={模型输出}]
\textbf{答}\enspace 答案是 24。\par
\hfill\textcolor{red!75!black}{\(\times\)}
\end{tcolorbox}
\end{minipage}\hfill
\begin{minipage}[t]{.502\linewidth}
\centering{\sffamily\bfseries 思维链提示}\par
\begin{tcolorbox}[cotpanel,title={模型输入}]
\textbf{问}\enspace Bob 举办聚会，买了 3 张披萨，每张切成 12 片。参加聚会的人吃掉了两整张披萨，以及另一张的三分之一。还剩多少片披萨？

\textbf{答}\enspace Bob 一开始有 3 张披萨，每张 12 片，也就是 \(12+12+12=36\) 片。吃掉两整张，还剩 \(36-12-12=12\) 片。最后一张又吃了三分之一，所以剩 \(12-4=8\) 片，因为 \(12/3=4\)。答案是 8。

\textbf{问}\enspace Rebecca 举办聚会，买了 4 盒甜甜圈，每盒 12 个。聚会上有 12 人，每人吃 2 个甜甜圈。还剩多少盒甜甜圈？
\end{tcolorbox}
\begin{tcolorbox}[cotpanel,title={模型输出}]
\textbf{答}\enspace 4 盒甜甜圈每盒 12 个，共有 \(12+12+12+12=48\) 个。12 人每人吃 2 个，还剩 \(48-24=24\) 个。每盒 12 个，所以还剩 2 盒甜甜圈。\par
\hfill\textcolor{green!50!black}{\bfseries 正确}
\end{tcolorbox}
\end{minipage}
\endgroup

\caption{假想的标准提示和思维链提示的比较。给大语言模型增加推理说明，既能提高答案的准确率，也使模型得出答案的过程更容易解释。}
\label{fig:ch01-cot}
\end{figure}

\subsubsection{应用}
如本章开头所言，公众最熟悉的大语言模型应用，是让模型写作或辅助写作。这类应用通常使用自回归模型。模型从输入文本开始，预测下一个词元，把这个词元追加到输入的末尾，然后以加长后的文本重复这个过程，直至完成生成。第~\ref{sec:intro-opening} 节提到的聊天机器人也以此为基础。人们还用这种方法开发辅助写作软件，帮助起草电子邮件、法律文书、技术手册、数据报告，以及几乎任何可以想到的文字材料。要为某个场景获得符合要求的高质量输出，通常需要用人工审定的数据，例如问答对，微调基础大语言模型，以引导应用给出所需的答案（见第 5 章）。

辅助写作在计算机编程领域也有显著应用。基本思想与自然语言写作相同，但训练数据换成了多种编程语言的大量代码。编程助手可以实时建议自动补全，可以根据伪代码生成与现有代码库风格相容的函数，也可以根据自然语言指令把数据库中的数据填入仪表板。GitHub Copilot 等应用已在程序员中广泛使用，既可简化软件开发中较烦琐的环节，也能增强他们的创造力和解决一般问题的能力。

对话、搜索或检索领域的另一些应用，尝试减轻大语言模型失效的负面影响，例如模型容易产生幻觉，即对查询给出不符合事实的回答。克服这一问题的一种方法是检索增强生成（RAG）。RAG 将一个大语言模型和一个特定主题的知识库配对，以便从中取得上下文。知识库可以是一份汽车使用手册，也可以是某一法律案件的一组文档。RAG 流程接收用户查询，判断查询是否与指定主题有关，然后搜索配对的知识库，提取与用户问题有关的信息。随后将知识库中相关的上下文与原始查询一同交给大语言模型，由它生成回答。这样，RAG 就能利用大语言模型准确回答与所给上下文有关的问题的能力。在推理时提供上下文，模型就无需事先包含那些信息。第 7 章将详细讨论 RAG 系统。

在重视隐私和数据保护的行业中，RAG 尤为常见。这些行业中的领域知识大多超出通用大语言模型的训练数据分布，因此开箱即用的对话应用价值可能有限。直接用受保护的信息，如患者病史，对模型作领域适配也不可取。用于训练大语言模型的信息可能深嵌在模型权重中，日后回答任何一次用户查询时都有泄露的风险。RAG 使受保护信息可以留在大语言模型的训练数据之外，同时仍属于模型可使用的知识范围。

在法律或自然语言处理研究等领域，大语言模型具有当时先进的实体识别与推理能力。过去，数据科学家团队要花几个月开发高质量的实体识别模型。他们采用耗时耗资源的全监督方法，并配合复杂的业务规则系统，才能围绕识别出的相关实体作准确推理并采取行动。如今，基于大语言模型的应用已经能识别这些实体并据此推理，其能力接近律师资格考试成绩位于前 10\% 的人类律师。不过，这类评估能在多大程度上衡量大语言模型相对于人类的能力，仍不清楚（\refcite{ch01-martinez}{Martínez, 2023}）。

本节只举出了众多应用中的一小部分，应用范围还在不断扩大。本书各处还会讨论更多应用，尤其是第 8 章。各领域的从业者继续探索将大语言模型用于自身工作时的收益与局限，与这些技术有关的创新也将不断加快。大语言模型在许多学习、评估和生成基准上表现突出，人们自然希望把这些能力用于解决企业和消费者的问题。研究人员通过微调或上下文学习等技术，让大语言模型适应各种领域，由此产生了许多发挥其新能力的应用。我们希望本书能成为了解大语言模型及其用法的入门教材和参考资料。

% PDF 55--58；原书第 24--27 页。
\begin{chapterreferences}
\refentry{ch01-askell}{Amanda Askell et al. A general language assistant as a laboratory for alignment, 2021.}
\refentry{ch01-bahdanau}{Dzmitry Bahdanau, Kyunghyun Cho, and Yoshua Bengio. Neural machine translation by jointly learning to align and translate. CoRR, abs/1409.0473, 2014.}
\refentry{ch01-bengio}{Yoshua Bengio, Réjean Ducharme, and Pascal Vincent. A neural probabilistic language model. In \textit{Proceedings of the 13th International Conference on Neural Information Processing Systems}, pages 893–899. MIT Press, 2000.}
\refentry{ch01-brown-1992}{Peter F. Brown, Peter V. deSouza, Robert L. Mercer, Vincent J. Della Pietra, and Jenifer C. Lai. Class-based n-gram models of natural language. \textit{Comput. Linguist.}, 18(4):467–479, December 1992.}
\refentry{ch01-brown-2020}{Tom Brown, Benjamin Mann, Nick Ryder, Melanie Subbiah, Jared D Kaplan, Prafulla Dhariwal, Arvind Neelakantan, Pranav Shyam, Girish Sastry, Amanda Askell, et al. Language models are few-shot learners. \textit{Advances in Neural Information Processing Systems}, 33:1877–1901, 2020.}
\refentry{ch01-chomsky}{Noam Chomsky. \textit{Syntactic Structures}. Mouton and Co., 1957.}
\refentry{ch01-chowdhery}{Aakanksha Chowdhery et al. PaLM: Scaling language modeling with Pathways, 2022.}
\refentry{ch01-collobert}{Ronan Collobert and Jason Weston. A unified architecture for natural language processing: Deep neural networks with multitask learning. In \textit{Proceedings of the 25th International Conference on Machine Learning}, pages 160–167. ACM, 2008.}
\refentry{ch01-devlin}{Jacob Devlin, Ming-Wei Chang, Kenton Lee, and Kristina Toutanova. BERT: Pre-training of deep bidirectional transformers for language understanding, 2019.}
\refentry{ch01-di-placido}{Dani Di Placido. Why did “Balenciaga Pope” go viral?, Mar 2023. URL \url{https://www.forbes.com/sites/danidiplacido/2023/03/27/why-did-balenciaga-pope-go-viral/}.}
\refentry{ch01-fortune}{Fortune. Generative AI market size, share COVID-19 impact analysis, by model, by industry vs application, and regional forecast, 2023–2030, Aug 2023. URL \url{https://www.fortunebusinessinsights.com/generative-ai-market-107837}.}
\refentry{ch01-gholami}{Amir Gholami, Sehoon Kim, Zhen Dong, Zhewei Yao, Michael W. Mahoney, and Kurt Keutzer. A survey of quantization methods for efficient neural network inference, 2021.}
\refentry{ch01-graves}{Alex Graves. Generating sequences with recurrent neural networks. CoRR, abs/1308.0850, 2013.}
\refentry{ch01-gururangan}{Suchin Gururangan, Ana Marasović, Swabha Swayamdipta, Kyle Lo, Iz Beltagy, Doug Downey, and Noah A. Smith. Don't stop pretraining: Adapt language models to domains and tasks, 2020.}
\refentry{ch01-hajicova}{Eva Hajicová, Ivana Kruijff-Korbayová, and Petr Sgall. Prague dependency treebank: Restoration of deletions. In \textit{Proceedings of the Second International Workshop on Text, Speech and Dialogue}, pages 44–49. Springer-Verlag, 1999.}
\refentry{ch01-hoffmann}{Jordan Hoffmann et al. Training compute-optimal large language models, 2022.}
\refentry{ch01-hu}{Edward J. Hu, Yelong Shen, Phillip Wallis, Zeyuan Allen-Zhu, Yuanzhi Li, Shean Wang, Lu Wang, and Weizhu Chen. LoRA: Low-rank adaptation of large language models, 2021.}
\refentry{ch01-hutchins}{W. John Hutchins, Leon Dostert, and Paul Garvin. The Georgetown-I.B.M. experiment. In In, pages 124–135. John Wiley And Sons, 1955.}
\refentry{ch01-ivison}{Hamish Ivison, Akshita Bhagia, Yizhong Wang, Hannaneh Hajishirzi, and Matthew Peters. HINT: Hypernetwork instruction tuning for efficient zero- few-shot generalisation, 2023.}
\refentry{ch01-jin}{Berber Jin and Miles Kruppa. WSJ News Exclusive \textbar{} ChatGPT creator is talking to investors about selling shares at \$29 billion valuation, Feb 2023. URL \url{https://www.wsj.com/articles/chatgpt-creator-openai-is-in-talks-for-tender-offer-that-would-value-it-at-29-billion-11672949279}.}
\refentry{ch01-kaufmann}{Timo Kaufmann, Paul Weng, Viktor Bengs, and Eyke Hüllermeier. A survey of reinforcement learning from human feedback, 2024.}
\refentry{ch01-lin}{Bill Yuchen Lin, Abhilasha Ravichander, Ximing Lu, Nouha Dziri, Melanie Sclar, Khyathi Chandu, Chandra Bhagavatula, and Yejin Choi. The unlocking spell on base LLMs: Rethinking alignment via in-context learning, 2023.}
\refentry{ch01-lv}{Kai Lv, Yuqing Yang, Tengxiao Liu, Qinghui Gao, Qipeng Guo, and Xipeng Qiu. Full parameter fine-tuning for large language models with limited resources, 2023.}
\refentry{ch01-ma}{Yingwei Ma, Yue Liu, Yue Yu, Yuanliang Zhang, Yu Jiang, Changjian Wang, and Shanshan Li. At which training stage does code data help LLMs reasoning?, 2023.}
\refentry{ch01-manning}{Christopher D. Manning and Hinrich Schütze. \textit{Foundations of Statistical Natural Language Processing}. MIT Press, 1999.}
\refentry{ch01-marcus}{Mitchell Marcus, Grace Kim, Mary Ann Marcinkiewicz, Robert MacIntyre, Ann Bies, Mark Ferguson, Karen Katz, and Britta Schasberger. The Penn Treebank: Annotating predicate argument structure. In \textit{Proceedings of the Workshop on Human Language Technology}, pages 114–119. Association for Computational Linguistics, 1994.}
\refentry{ch01-martinez}{Eric Martínez. Re-evaluating GPT-4’s bar exam performance. 2023. URL \url{http://dx.doi.org/10.2139/ssrn.4441311}.}
\refentry{ch01-mikolov-2010}{Tomas Mikolov, Martin Karafiát, Lukás Burget, Jan Cernocký, and Sanjeev Khundanpur. Recurrent neural network based language model. In Takao Kobayashi, Keikichi Hirose, and Satoshi Nakamura, editors, \textit{INTERSPEECH}, pages 1045–1048. ISCA, 2010.}
\refentry{ch01-mikolov-2013a}{Tomas Mikolov, Kai Chen, Greg Corrado, and Jeffrey Dean. Efficient estimation of word representations in vector space. CoRR, abs/1301.3781, 2013a.}
\refentry{ch01-mikolov-2013b}{Tomas Mikolov, Ilya Sutskever, Kai Chen, Greg S Corrado, and Jeff Dean. Distributed representations of words and phrases and their compositionality. In C. J. C. Burges, L. Bottou, M. Welling, Z. Ghahramani, and K. Q. Weinberger, editors, \textit{Advances in Neural Information Processing Systems 26}, pages 3111–3119. Curran Associates, Inc., 2013b.}
\refentry{ch01-miller}{George A. Miller. WordNet: A lexical database for English. \textit{Commun. ACM}, 38(11):39–41, November 1995.}
\refentry{ch01-openai}{OpenAI. GPT-4 technical report, 2023.}
\refentry{ch01-ouyang}{Long Ouyang et al. Training language models to follow instructions with human feedback, 2022.}
\refentry{ch01-radford}{Alec Radford, Jeffrey Wu, Rewon Child, David Luan, Dario Amodei, Ilya Sutskever, et al. Language models are unsupervised multitask learners. \textit{OpenAI blog}, 1(8):9, 2019.}
\refentry{ch01-rae}{Jack W. Rae et al. Scaling language models: Methods, analysis insights from training Gopher, 2022.}
\refentry{ch01-rafailov}{Rafael Rafailov, Archit Sharma, Eric Mitchell, Stefano Ermon, Christopher D. Manning, and Chelsea Finn. Direct preference optimization: Your language model is secretly a reward model, 2023.}
\refentry{ch01-raffel}{Colin Raffel, Noam Shazeer, Adam Roberts, Katherine Lee, Sharan Narang, Michael Matena, Yanqi Zhou, Wei Li, and Peter J. Liu. Exploring the limits of transfer learning with a unified text-to-text Transformer, 2020.}
\refentry{ch01-soares}{Livio Baldini Soares, Nicholas FitzGerald, Jeffrey Ling, and Tom Kwiatkowski. Matching the blanks: Distributional similarity for relation learning, 2019.}
\refentry{ch01-srivastava}{Aarohi Srivastava et al. Beyond the imitation game: Quantifying and extrapolating the capabilities of language models, 2023.}
\refentry{ch01-thoppilan}{Romal Thoppilan, Daniel De Freitas, Jamie Hall, Noam Shazeer, Apoorv Kulshreshtha, Heng-Tze Cheng, Alicia Jin, Taylor Bos, Leslie Baker, Yu Du, et al. LaMDA: Language models for dialog applications. \textit{arXiv preprint arXiv:2201.08239}, 2022.}
\refentry{ch01-touvron}{Hugo Touvron et al. Llama 2: Open foundation and fine-tuned chat models, 2023.}
\refentry{ch01-valuates}{Valuates. Generative AI market size to grow USD 126.5 billion by 2031 at a CAGR of 32\%, Jun 2023. URL \url{https://www.prnewswire.com/news-releases/generative-ai-market-size-to-grow-usd-126-5-billion-by-2031-at-a-cagr-of-32--valuates-reports-301846316.html}.}
\refentry{ch01-vaswani}{Ashish Vaswani, Noam Shazeer, Niki Parmar, Jakob Uszkoreit, Llion Jones, Aidan N Gomez, Łukasz Kaiser, and Illia Polosukhin. Attention is all you need. In I. Guyon, U. Von Luxburg, S. Bengio, H. Wallach, R. Fergus, S. Vishwanathan, and R. Garnett, editors, \textit{Advances in Neural Information Processing Systems}, volume 30. Curran Associates, Inc., 2017. URL \url{https://proceedings.neurips.cc/paper_files/paper/2017/file/3f5ee243547dee91fbd053c1c4a845aa-Paper.pdf}.}
\refentry{ch01-wei-2023}{Jason Wei, Xuezhi Wang, Dale Schuurmans, Maarten Bosma, Brian Ichter, Fei Xia, Ed Chi, Quoc Le, and Denny Zhou. Chain-of-thought prompting elicits reasoning in large language models, 2023.}
\refentry{ch01-wei-2022}{Jason Wei et al. Emergent abilities of large language models, 2022.}
\refentry{ch01-wu}{Shijie Wu, Ozan Irsoy, Steven Lu, Vadim Dabravolski, Mark Dredze, Sebastian Gehrmann, Prabhanjan Kambadur, David Rosenberg, and Gideon Mann. BloombergGPT: A large language model for finance, 2023.}
\refentry{ch01-zhang-a}{Qiang Zhang et al. Scientific large language models: A survey on biological chemical domains, 2024a.}
\refentry{ch01-zhang-b}{Shengyu Zhang et al. Instruction tuning for large language models: A survey, 2024b.}
\refentry{ch01-zhao}{Wayne Xin Zhao et al. A survey of large language models, 2023.}
\end{chapterreferences}

